We are now ready to define, for each $\dR\in \R$, a map $\Jhom_{\dR}$ which will eventually lead to $\hJhom_\dR$.

\begin{Definition}\label{D:theta}
Let $\lambda$ be a partition of $n$. Recall that there is a natural left $\sym{n}$-action on the set $\TT(\lambda)$ of $\lambda$-tableaux, and let $\Z\TT(\lambda)$ denote the associated permutation $\Z\sym{n}$-module. Let $(\alpha|\beta)$ be a bicomposition of $n$ and let $\Gamma$ be a fixed left transversal of $\sym{\alpha|\beta}$ in $\sym{n}$. For each $\dR \in\R$, we define a $\Z$-linear map $\Jhom_{\dR}\colon \Z\TT(\lambda) \to M_\Z(\alpha|\beta)$ as follows
\begin{gather*}\Jhom_{\dR}(\t) = \sum_{d \in \Gamma} \a_{\rho_{\t}^{-1} d,\dR} (d \otimes \bb{1} \otimes \epsilon),\end{gather*}
where $\rho_{\t} = \t \circ (\t_0)^{-1} \in \sym{n}$ (so that $\rho_{\t} \cdot \t_0 = \rho_{\t} \circ \t_0 = \t$).
\end{Definition}

The following properties of the map $\Jhom_{\dR}$ follow easily from its definition and Lem\-ma~\ref{L:a}(i).

\begin{Lemma} \label{L:Jhom}
Let $\dR \in \R$.
\begin{enumerate}\itemsep=0pt
 \item [$(i)$] The map $\Jhom_{\dR}$ is independent of the choice of the left transversal $\Gamma$, i.e., if $\Gamma'$ is any left transversal of $\sym{\alpha|\beta}$ in $\sym{n}$, then
 \begin{gather*}
 \Jhom_{\dR}(\t) = \sum_{d' \in \Gamma'} \a_{\rho_{\t}^{-1} d',\dR} (d' \otimes \bb{1} \otimes \epsilon)\end{gather*} for all $\t \in \TT(\lambda)$.
 \item [$(ii)$] If $\dR' \in R_{\t_0} \dR \sym{\alpha|\beta}$, then $\Jhom_{\dR'} = \e_{\dR}(\dR') \Jhom_{\dR}$.
\end{enumerate}
\end{Lemma}

\begin{proof} Let $d' = d \eta_{\alpha} \eta_{\beta}^{+|\alpha|}$ with $(\eta_{\alpha},\eta_{\beta}) \in \sym{\alpha} \times \sym{\beta}$. Then $\a_{\rho_{\t}^{-1}d',\dR} = \sgn(\eta_\beta) \a_{\rho_{\t}^{-1}d,\dR}$ by Lem\-ma~\ref{L:a}(i) while
\begin{gather*}d' \otimes \bb{1} \otimes \epsilon = (d \eta_{\alpha} \eta_{\beta}) \otimes \bb{1} \otimes \epsilon = d \otimes (\eta_{\alpha} \cdot \bb{1}) \otimes (\eta_{\beta} \cdot \epsilon) = \sgn(\eta_{\beta}) (d \otimes \bb{1} \otimes \epsilon).\end{gather*}
Thus $\a_{\rho_{\t}^{-1}d',\dR} (d' \otimes \bb{1} \otimes \epsilon) = \a_{\rho_{\t}^{-1}d,\dR} (d \otimes \bb{1} \otimes \epsilon)$ and part (i) follows.

By Lemma \ref{L:a}(i), $\a_{\rho_{\t}^{-1}d,\dR'} = \e_{\dR}(\dR') \a_{\rho_{\t}^{-1}d,\dR}$. Part (ii) thus follows.
\end{proof}

Next, we aim to show that each of the maps $\Jhom_{\dR}\colon \Z\TT(\lambda)\to M_\Z(\alpha|\beta)$ induces a $\Z\sym{n}$-module homomorphism $\hJhom_{\dR}\colon S^\lambda_\Z\to M_\Z(\alpha|\beta)$. For this, we will show that $\Jhom_{\dR}$ is a $\Z\sym{n}$-module homomorphism and that the kernel of the natural map $\psi\colon \Z\TT(\lambda)\to S^\lambda_\Z$ given by $\psi(\t)=e_\t$ is contained in $\ker(\Jhom_{\dR})$ and hence $\Jhom_{\dR}$ induces $\hJhom_{\dR}$ as desired.

For each $j$, let $C_j(\lambda) = \{ (i,j) \in [\lambda] \colon 1 \leq i \leq \ell(\lambda) \}$ be the $j$th column of the Young diagram of $\lambda$. A {\em Garnir transversal} $\Delta$ is a left transversal of $\sym{X}\sym{Y}$ in $\sym{X \cup Y}$, where $\varnothing \ne X \subseteq C_j(\lambda)$ and $\varnothing \ne Y \subseteq C_{j'}(\lambda)$ with $j < j'$, such that $|X|+|Y| > |C_j(\lambda)|$. Given a Garnir transversal $\Delta$ and a $\lambda$-tableau $\t$, let $\Delta_{\t} = \{ \t \circ \gamma \circ \t^{-1} \colon \gamma \in \Delta \}$, so that $\Delta_{\t}$ is a left transversal of $\sym{\t(X)}\sym{\t(Y)}$ in~$\sym{\t(X \cup Y)}$, and write
\begin{gather*}G_{\Delta}^{\t} = \sum_{\gamma \in \Delta_\t} \sgn(\gamma) \gamma.\end{gather*}

\begin{Proposition}\label{P: phi_d hom}
Let $\dR \in \R$. Then
\begin{enumerate}\itemsep=0pt
\item [$(i)$] $\Jhom_{\dR}$ is a $\Z \sym{n}$-module homomorphism, and
\item [$(ii)$] for any $\t \in \TT(\lambda)$, $\pi\in C_\t$ and Garnir transversal $\Delta$, we have that $\ker(\Jhom_{\dR})$ contains both $\pi \cdot \t - \sgn(\pi) \t$ and $G_{\Delta}^{\t} \cdot \t$.
\end{enumerate}
\end{Proposition}

\begin{proof} Let $\t \in \TT(\lambda)$ and $x \in \sym{n}$. We have
\begin{gather*}
\Jhom_{\dR}(x \cdot \t) = \Jhom_{\dR}((x \rho_{\t}) \cdot \t_0) = \sum_{d \in \Gamma} \a_{(x \rho_{\t})^{-1} d, \dR} (d \otimes \bb{1} \otimes \epsilon), \\
x \cdot \Jhom_{\dR} (\t) = x \cdot \left (\sum_{d \in \Gamma} \a_{\rho_{\t}^{-1} d, \dR} (d \otimes \bb{1} \otimes \epsilon)\right ) = \sum_{d \in \Gamma} \a_{\rho_{\t}^{-1} d, \dR} \sgn(\xi_{d,\beta}) (f(d) \otimes \bb{1} \otimes \epsilon),
\end{gather*}
where $xd = f(d) \xi_{d,\alpha} \xi_{d,\beta}^{+|\alpha|}$ with $f(d) \in \Gamma$, $(\xi_{d,\alpha},\xi_{d,\beta}) \in \sym{\alpha} \times \sym{\beta}$. As such, we need to show that for each $d \in \Gamma$,
\begin{gather*}
\a_{(x \rho_{\t})^{-1} f(d), \dR} = \a_{\rho_{\t}^{-1} d, \dR} \sgn(\xi_{d,\beta}).
\end{gather*}
Observe that $\sigma (x\rho_{\t})^{-1}f(d)=\tau \dR \eta_\alpha \eta_\beta^{+|\alpha|}$ for some $\sigma\in C_{\t_0}$, $\tau\in R_{\t_0}$ and $(\eta_\alpha,\eta_\beta)\in\sym{\alpha}\times \sym{\beta}$ if and only if $\sigma\rho_{\t}^{-1}d=\tau \dR (\eta_\alpha \xi_{d,\alpha}) (\eta_\beta\xi_{d,\beta})^{+|\alpha|}$ for some $\sigma\in C_{\t_0}$, $\tau\in R_{\t_0}$ and $(\eta_\alpha,\eta_\beta)\in\sym{\alpha}\times \sym{\beta}$, and in which case, $\e_{\dR}\big( \sigma (\rho_{\t})^{-1}d\big) = \sgn(\xi_{d,\beta}) \e_{\dR}\big(\sigma(x\rho_{\t})^{-1}f(d)\big)$. So $\Omega_{(x\rho_{\t})^{-1}f(d),\dR}=\Omega_{\rho_{\t}^{-1}d,\dR}$, and
\begin{align*}
\a_{(x\rho_{\t})^{-1} f(d),\dR} &= \sum_{\sigma\in \Omega_{(x\rho_{\t})^{-1}f(d),\dR}} \sgn(\sigma)\e_{\dR}\big(\sigma (x\rho_{\t})^{-1} f(d)\big)\\
&=\sum_{\sigma\in \Omega_{\rho_{\t}^{-1}d,\dR}} \sgn(\sigma)\sgn(\xi_{d,\beta})\e_{\dR}\big(\sigma \rho_{\t}^{-1}d\big)
=\a_{\rho_{\t}^{-1}d,\dR} \sgn(\xi_{d,\beta}).
\end{align*} This proves part (i).

Since $\rho_\t\cdot \t_0=\t$, we have $C_\t=\rho_\t C_{\t_0}\rho_\t^{-1}$ and hence $\rho_\t^{-1}\pi^{-1}\rho_\t\in C_{\t_0}$. By Lemma \ref{L:a}(i), we have \begin{gather*}\a_{(\pi\rho_{\t})^{-1} d, \dR}=\a_{\rho_\t^{-1}\pi^{-1}\rho_\t\rho^{-1}_\t d, \dR} = \sgn(\pi) \a_{\rho_{\t}^{-1}d, \dR}.\end{gather*} Therefore
\begin{gather*}
\Jhom_{\dR}(\pi \cdot \t) = \sum_{d \in \Gamma} \a_{(\pi \rho_{\t})^{-1}d, \dR} (d \otimes \bb{1} \otimes \epsilon)
=\sum_{d \in \Gamma} \sgn(\pi) \a_{\rho_{\t}^{-1}d, \dR} (d \otimes \bb{1} \otimes \epsilon)
=\sgn(\pi)\Jhom_{\dR}(\t).
\end{gather*}

Next we turn to $G_{\Delta}^{\t} \cdot \t$. We have
\begin{align*}
\Jhom_{\dR}\big(G_{\Delta}^{\t} \cdot \t\big) & = \Jhom_{\dR}\left (\sum_{\gamma \in \Delta_{\t}} \sgn(\gamma) (\gamma \cdot \t)\right )
= \sum_{\gamma \in \Delta_{\t}} \sum_{d \in \Gamma} \sgn(\gamma) \a_{(\gamma \rho_{\t})^{-1} d, \dR} (d \otimes \bb{1} \otimes \epsilon) \\
& = \sum_{d \in \Gamma} \left( \sum_{\gamma \in \Delta_{\t}} \sum_{\sigma \in \Omega_{(\gamma \rho_{\t})^{-1} d, \dR}} \sgn(\gamma) \sgn(\sigma) \e_{\dR}\big( \sigma (\gamma \rho_{\t})^{-1} d\big) \right) (d \otimes \bb{1} \otimes \epsilon).
\end{align*}
Fix $d \in \Gamma$, and let $b(\gamma,\sigma) = \sgn(\gamma) \sgn(\sigma) \e_{\dR}\big( \sigma (\gamma \rho_{\t})^{-1} d\big)$. We need to show that $\!\sum\limits_{(\gamma,\sigma) \in \Upsilon} \! b(\gamma,\sigma) =0$, where $\Upsilon = \{ (\gamma,\sigma) \colon \gamma \in \Delta_{\t},\, \sigma \in \Omega_{(\gamma \rho_{\t})^{-1} d, \dR} \}$.

Let the Garnir transversal $\Delta$ be a left transversal of $\sym{X}\sym{Y}$ in $\sym{X \cup Y}$, where $X$ and $Y$ are subsets of the $j$th and $j'$th column of $[\lambda]$, with $|X| + |Y| > |C_j(\lambda)|$. Since $|X| + |Y| > |C_j(\lambda)|$, there exists $1 \leq i \leq \ell(\lambda)$ such that $(i,j) \in X$ and $(i,j') \in Y$, and we may choose $i$ to be the least such. Let $\eta$ be the transposition $(\t_0(i,j) \;\; \t_0(i,j'))$. Then $\eta \in R_{\t_0}$.

Let $(\gamma, \sigma) \in \Upsilon$, say $\sigma (\gamma \rho_{\t})^{-1} d = \tau \dR \xi_\alpha \tlxi{\beta}{|\alpha|}$ where $\tau \in R_{\t_0}$ and $(\xi_\alpha,\xi_\beta) \in \sym{\alpha} \times \sym{\beta}$. Then
\begin{align*}
R_{\t_0} \dR \sym{\alpha|\beta} \ni \eta \tau \dR \xi_\alpha \tlxi{\beta}{|\alpha|} &=
\eta \sigma (\gamma \rho_{\t})^{-1} d = \eta \sigma \rho_{\t}^{-1} \gamma^{-1} d\\
&= \big(\sigma \rho_{\t}^{-1}\big) \big(\big(\sigma \rho_{\t}^{-1}\big)^{-1} \eta \big(\sigma \rho^{-1}_{\t} \big)\big) \gamma^{-1} d
= \sigma \rho_{\t}^{-1} \zeta^{-1} \rho_{\t} (\gamma'\rho_{\t})^{-1} d,
\end{align*}
where $\gamma' \in \Delta_{\t}$ and $\zeta \in \sym{\t(X)}\sym{\t(Y)}$ satisfies $\gamma \big(\big(\sigma \rho_{\t}^{-1}\big)^{-1} \eta \big(\sigma \rho^{-1}_{\t} \big)\big) = \gamma' \zeta$. Then $\rho_{\t}^{-1} \zeta^{-1} \rho_{\t} \in \rho_{\t}^{-1} \sym{\t(X)} \sym{\t(Y)} \rho_{\t} = \sym{\t_0(X)}\sym{\t_0(Y)} \subseteq C_{\t_0}$. This shows that $(\gamma',\sigma \rho_{\t}^{-1} \zeta^{-1} \rho_{\t}) \in \Upsilon$. In other words, the function $h \colon \Upsilon \to \Upsilon$ defined by $(\gamma,\sigma) \mapsto \big(\gamma',\sigma \rho_{\t}^{-1} \zeta^{-1} \rho_{\t}\big)$, where $\gamma \big(\big(\sigma \rho_{\t}^{-1}\big)^{-1} \eta \big(\sigma \rho^{-1}_{\t} \big)\big) = \gamma' \zeta$, is well-defined. Furthermore,
\begin{align*}
b(h(\gamma,\sigma)) &= \sgn(\gamma')\sgn\big(\sigma \rho_{\t}^{-1} \zeta^{-1} \rho_{\t} \big) \e_{\dR}\big(\sigma \rho_{\t}^{-1} \zeta^{-1} \rho_{\t} (\gamma'\rho_{\t})^{-1} d\big) \\
&= -\sgn(\gamma) \sgn(\sigma) \sgn(\xi_{\beta}) = -b(\gamma,\sigma).
\end{align*}
We claim that $h$ is a fixed-point-free involution, in which case, since the contributions from $(\gamma,\sigma)$ and $h(\gamma,\sigma)$ towards the sum $\sum\limits_{(\gamma,\sigma) \in \Upsilon} b(\gamma,\sigma)$ cancel each other out, we have $\sum\limits_{(\gamma,\sigma) \in \Upsilon} b(\gamma,\sigma) =0$ as desired.

To prove the claim, first suppose that
\begin{gather*}
(\gamma,\sigma) =h(\gamma,\sigma)=\big(\gamma',\sigma \rho_{\t}^{-1} \zeta^{-1} \rho_{\t}\big).\end{gather*}
Then $\gamma=\gamma'$ and $\zeta=1$ and hence $\eta=1$, a contradiction. So $h$ is fixed-point-free. Next, $h^2(\gamma,\sigma) = h\big(\gamma',\sigma \rho_{\t}^{-1} \zeta^{-1} \rho_{\t}\big) = \big(\gamma'', \sigma \rho_{\t}^{-1} \zeta^{-1} \rho_{\t} \rho_{\t}^{-1} \zeta'^{-1} \rho_{\t}\big)$, where
\begin{gather*}
\gamma''\zeta' = \gamma' \big(\sigma \rho_{\t}^{-1} \zeta^{-1}\big)^{-1} \eta \big(\sigma \rho_{\t}^{-1} \zeta^{-1}\big) =
\gamma' \zeta \big(\sigma \rho_{\t}^{-1}\big)^{-1} \eta \big(\sigma \rho_{\t}^{-1}\big) \zeta^{-1} = \gamma \zeta^{-1}.
\end{gather*}
Thus $\gamma'' = \gamma$ and $\zeta' = \zeta^{-1}$ and hence $h^2(\gamma,\sigma) = (\gamma, \sigma)$, and the proof is complete.
\end{proof}

The next result is well known when the underlying ring is a field (see, for example, \cite[Section~7.4, Corollary, p.~101]{WF}); we are however unable to find its generalisation to~$\mathbb{Z}$ in the existing literature.

\begin{Lemma}\label{L: annihilator for specht}
The map $\psi\colon \Z\TT(\lambda) \to S_\Z^\lambda$ defined by $\t \mapsto e_{\t}$ is a $\Z\sym{n}$-module epimorphism, and $\ker(\psi)$ is generated, as a $\Z$-submodule of $\Z\TT(\lambda)$, by $G \cup H$, where
\begin{gather*}
G= \big\{ G_{\Delta}^{\t} \cdot \t \colon \t \in \TT(\lambda),\ \Delta \text{ a Garnir transversal} \big\}, \\
H= \{ \pi \cdot \t - \sgn(\pi)\t \colon \t \in \TT(\lambda),\ \pi \in C_{\t} \}.
\end{gather*}
\end{Lemma}

\begin{proof}That $\psi$ is a $\Z\sym{n}$-module epimorphism is clear, so we only need to justify the assertion about its kernel. Let $K = \Z(G\cup H)$. It is straightforward to verify that both $G$ and $H$ are invariant under the action of $\sym{n}$, and that $\psi(G \cup H) = \{0\}$ \cite[7.2.1 and Theorem~7.2.3]{GJAK}, so that~$K$ is a $\Z\sym{n}$-submodule of $\Z\TT(\lambda)$, and $K \subseteq \ker(\psi)$. It remains to show that $\ker(\psi) \subseteq K$.

Since $\psi(\Z\TT(\lambda)) = S^{\lambda}_\Z$, and $S_\Z^{\lambda}$ is $\Z$-free with basis $\{ e_{\t} \colon \t \in \TT_{\mathrm{std}}(\lambda) \}$ \cite[Theorem~1.1]{Peel}, we see that $\Z\TT(\lambda) = (\bigoplus_{\t \in \TT_{\mathrm{std}}(\lambda)} \Z\t)\oplus \ker(\psi)$.

Consider $\Z\TT(\lambda)/K$. For each $\t \in \TT(\lambda)$, let $v_{\t} = \t + K \in \Z\TT(\lambda)/K$. Then $\pi\cdot v_\t = \sgn(\pi)v_\t$ for any $\pi \in C_{\t}$, and $G_{\Delta}^{\t} \cdot v_\t = 0$ for any Garnir transversal $\Delta$. Using the same argument as in the proof of \cite[Theorem~7.2.7]{GJAK}, we can show that $\{ v_{\t} \colon \t \in \TT_{\mathrm{std}}(\lambda) \}$ generates $\Z\TT(\lambda)/K$ as a $\Z$-module.

Let $x \in \ker(\psi)$. Then there exists $b_{\t} \in \Z$ for each $\t \in \TT_{\mathrm{std}}(\lambda)$ such that \begin{gather*}x+K = \sum_{\t \in \TT_{\mathrm{std}}(\lambda)} b_{\t}v_{\t} = \sum_{\t \in \TT_{\mathrm{std}}(\lambda)} b_{\t}(\t + K).\end{gather*}
Thus there exists $k \in K$ such that \begin{gather*}x + k = \sum_{\t \in \TT_{\mathrm{std}}(\lambda)} b_{\t}\t \in \left (\bigoplus_{\t \in \TT_{\mathrm{std}}(\lambda)} \Z\t\right ) \cap \ker(\psi) = \{0\}.\end{gather*} Hence $x + k = 0$ and so $x \in K$, and our proof is complete.
\end{proof}

We are now ready to prove Theorem \ref{T: hom}.

\begin{proof}[Proof of Theorem \ref{T: hom}]
By Proposition \ref{P: phi_d hom} and Lemma \ref{L: annihilator for specht}, we have $\ker(\psi) \subseteq \ker(\Jhom_{\dR})$, so that there is a unique $\Z\sym{n}$-module homomorphism $\hJhom_{\dR}\colon S_\Z^\lambda \to M_\Z(\alpha|\beta)$ such that $\hJhom_{\dR} \circ \psi = \Jhom_{\dR}$. The theorem immediately follows.
\end{proof}

Let $\F$ be a field. Since
\begin{gather*}
S_\F^{\lambda} \cong \F \otimes_{\mathbb{Z}} S_\Z^{\lambda} \qquad \text{and} \qquad M_\F(\alpha|\beta) \cong \F \otimes_{\mathbb{Z}} M_\Z(\alpha|\beta),
\end{gather*}
each $\hJhom_{\dR}$ gives rise to an $\F\sym{n}$-homomorphism $\hJhom_{\dR}^{\F} \colon S_\F^{\lambda} \to M_\F(\alpha|\beta)$. More specifically, for any $k \in \mathbb{Z}$, let $k^{\F} =k \cdot 1_{\F} \in \F$. Then
\begin{gather*}\hJhom_{\dR}^{\F} (e_{\t}) = \sum_{d\in\Gamma} \a_{\rho_{\t}^{-1}d,\dR}^{\F} \,(d \otimes \bb{1} \otimes \epsilon).\end{gather*}

As we mentioned in the paragraph following Theorem \ref{T: hom}, our $\hJhom_{\dR}^{\F}$'s generalise the $\hat{\theta}_{\T}$'s constructed by James.
